\documentclass[11pt]{article}
\usepackage[T1]{fontenc}
\usepackage[margin=1in]{geometry}
\usepackage{amsmath,amssymb}
\usepackage{enumitem}
\usepackage{hyperref}
\hypersetup{colorlinks=true,linkcolor=black,urlcolor=blue}

\begin{document}

\begin{center}
{\LARGE Breathing From ACC Formula Sheet}\\[0.5em]
{\large PolarH10 repository-specific breathing-volume approximation}
\end{center}

\section*{Provenance}

This breathing-volume model is repository-specific code ported from the existing
Unity runtime. It is not presented as a published external method reference.

\begin{itemize}[leftmargin=1.5em]
\item This ACC breathing-volume model has not yet been externally validated.
\item A validation study is actively being worked on.
\end{itemize}

\section*{Data path}

\begin{itemize}[leftmargin=1.5em]
\item input stream: PMD accelerometer samples from the Polar H10
\item preprocessing: exponential moving average on the raw ACC signal
\item calibration: estimate a center point and a principal breathing axis from a calibration window
\item projection: map filtered ACC onto the chosen breathing axis
\item bounds: derive lower and upper projection bounds from calibration quantiles
\item output: convert the projected signal into a $0..1$ breathing-volume estimate
\end{itemize}

\section*{Core approximation}

Filtered ACC sample:

\[
a_f(t) = (1 - \alpha)a_f(t-1) + \alpha a_{raw}(t)
\]

Calibration center:

\[
c = \frac{1}{N}\sum_{i=1}^{N} a_f(i)
\]

Principal-axis projection:

\[
p(t) = \langle a_f(t) - c, u \rangle
\]

where $u$ is the dominant calibration axis estimated from the calibration-window
covariance.

If the optional XZ model is available, the tracker also keeps a 2-D projection
in the chest plane:

\[
p_{xz}(t) = \langle (x(t), z(t)) - c_{xz}, u_{xz} \rangle
\]

Lower and upper projection bounds come from calibration quantiles:

\[
b_{min} = Q_{lower}(p), \qquad b_{max} = Q_{upper}(p)
\]

The runtime then applies the configured edge-ease and optional adaptive-bounds
updates before mapping the projection to volume:

\[
v(t) = \operatorname{clamp}\left(\frac{p(t) - b_{min}}{b_{max} - b_{min}}, 0, 1\right)
\]

The final base volume uses either the 3-D axis model or the XZ model, depending
on the tracker settings. The displayed output may then be inverted if the
configured inhale/exhale direction is flipped.

\section*{State classification}

The tracker derives \texttt{Inhaling}, \texttt{Exhaling}, and \texttt{Pausing}
from the first difference of the output volume:

\[
\Delta v(t) = v(t) - v(t-1)
\]

\begin{itemize}[leftmargin=1.5em]
\item \texttt{Inhaling} when $\Delta v$ exceeds the configured positive threshold
\item \texttt{Exhaling} when $\Delta v$ is below the negative threshold
\item \texttt{Pausing} otherwise
\end{itemize}

\section*{Alignment status}

There is no external paper claim here. The correct interpretation is:

\begin{itemize}[leftmargin=1.5em]
\item this sheet documents what the app computes
\item the model is intended as a practical ACC-only breathing approximation
\item external validation is still pending
\end{itemize}

\section*{References}

\begin{itemize}[leftmargin=1.5em]
\item \url{https://mesmerprism.github.io/PolarH10/reference/breathing-workflow.html}
\item \url{https://mesmerprism.github.io/PolarH10/reference/references.html}
\end{itemize}

\end{document}
